%!TeX root=MemoriaTFG.tex

\chapter{Introducción}
\label{cap:intro}

\section{Contexto clínico}
\label{sec:contexto}

Aunque existen centenares de enfermedades cutáneas distintas, la
actividad asistencial diaria de los servicios de dermatología se
concentra principalmente en un número limitado de diagnósticos
frecuentes. El estudio observacional DIADERM, sobre $10\,999$
diagnósticos en $8\,953$ pacientes en territorio español, sitúa
entre los diagnósticos más prevalentes la queratosis actínica, el
carcinoma basocelular, los nevus melanocíticos, la queratosis
seborreica, la psoriasis, el acné y las verrugas
víricas~\cite{diaderm2018}. Por su parte, Taberner et~al.\
documentaron $213$ diagnósticos diferentes en un único año
asistencial en el Hospital Universitario Son Llàtzer
(HUSLL)~\cite{taberner2010_motivos}. Ambos trabajos reflejan una
realidad común: una elevada diversidad diagnóstica sustentada sobre
un núcleo relativamente reducido de dermatosis benignas,
inflamatorias e infecciosas de alta frecuencia.

La sensibilidad temporal del proceso diagnóstico resulta crítica en
el melanoma cutáneo: la supervivencia específica a cinco años
desciende desde aproximadamente $99\%$ en estadio~I hasta $25\%$ en
estadio~IV según la octava edición de la American Joint Committee on
Cancer~\cite{ajcc8}. El acceso al especialista se realiza mediante
derivación desde atención primaria, y la capacidad del circuito
presencial es limitada frente a una demanda creciente, lo que motiva
el interés por sistemas de apoyo al diagnóstico basados en imagen. El
sustrato visual de ese diagnóstico se articula en tres modalidades
---fotografía clínica, dermatoscopia e histopatología---, sobre las
que operan hoy sistemas construidos a partir de \emph{modelos
fundacionales}: modelos de gran escala, preentrenados sobre un
volumen amplio de datos y reutilizables, mediante adaptación, en
múltiples tareas posteriores.

Este trabajo nace de dos colaboraciones preexistentes con la
Dra.~R.~Taberner, dermatóloga del HUSLL. La primera es la
modernización tecnológica del archivo Dermapixel, blog mantenido por
la Dra.~Taberner desde 2010, que reúne más de $700$ casos clínicos
comentados con imagen clínica y dermatoscópica, historia abreviada,
diagnóstico y razonamiento diferencial; en paralelo a este TFG se ha
desarrollado su migración a una pila tecnológica propia. La segunda
es un proyecto institucional del Servei de Salut de les Illes Balears
(IB-Salut) sobre teledermatología hospitalaria, orientado a apoyar la
primera valoración dermatológica en el circuito de derivación desde
atención primaria. \label{sec:trabajo-previo}Ambas líneas orientan
una decisión que condiciona todo el trabajo: priorizar modelos con
código y pesos públicos, por su viabilidad de integración
hospitalaria y su reproducibilidad. Esta prioridad introduce un sesgo
de selección explícito ---deja fuera modelos cerrados potencialmente
más capaces, como DermFM-Zero, cuyos pesos no son públicos
(sección~\ref{sec:lim-pesos})---, asumido a cambio de un marco de
evaluación íntegramente reproducible; los modelos cerrados accesibles
vía API se incorporan únicamente como comparadores contextuales.

\section{Objetivo general y enfoque}
\label{sec:objetivos}

El objetivo último de este trabajo es avanzar hacia un sistema de
apoyo diagnóstico capaz de operar tanto sobre fotografía clínica como
sobre dermatoscopia, y sobre el conjunto amplio de patologías
cutáneas, no solo el melanoma. La motivación es una asimetría clara
del estado actual: la dermatoscopia está relativamente madura para la
detección de melanoma, pero mucho menos para el resto de patologías,
mientras que la fotografía clínica ---la modalidad más accesible en
atención primaria y en teledermatología--- sigue muy poco explorada.
Ahí residen la necesidad clínica y la oportunidad.

El modo de proceder es, por tanto, empezar por el principio: examinar
qué ofrecen y qué les falta a los modelos fundacionales disponibles.
De ese análisis del estado del arte se derivarán los objetivos
concretos y el enfoque experimental de este trabajo.
